% ============================================================
\chapter{Background and Related Work}
\label{chap:background}
% ============================================================

Rather than survey the two literatures this project sits between --- portfolio
construction and causal discovery --- this chapter analyses in depth the three
works it directly builds on, one per theme. L\'opez de Prado's Hierarchical
Risk Parity and backtest-rigour agenda (Section~\ref{sec:hrp-theme}) supplies
both the allocator whose symmetric clustering step C2 targets and the
inferential standard C3 adopts. Howard, Lohre and Mudde's causal networks on
the S\&P~500 (Section~\ref{sec:howard-theme}) are the closest causal-finance
work, and stop exactly one step short of the question Q asks. Rodr\'iguez-%
Dom\'inguez's Hierarchical Sensitivity Parity (Section~\ref{sec:hsp-theme})
is the framework this project began from, and the instabilities identified
here are the ones the evaluation quantifies. Section~\ref{sec:cd-tools}
compresses the causal-discovery toolkit into the facts the pipeline relies
on, and Section~\ref{sec:gap} states the gap.

\section{Theme I --- hierarchical allocation and honest backtests:
L\'opez de Prado}
\label{sec:hrp-theme}

Mean--variance optimisation \cite{markowitz1952portfolio} requires inverting a
covariance matrix that, for a 100-asset universe estimated on a one-year
window, is noisy and near-singular; the optimiser amplifies the estimation
error and produces concentrated, unstable weights. Hierarchical Risk Parity
\cite{lopezdeprado2016hrp} avoids the inversion with a three-step scheme.
\textbf{(i) Clustering:} convert the correlation matrix $\rho$ into the
distance $d_{ij} = \sqrt{\tfrac{1}{2}(1-\rho_{ij})}$ and build a
single-linkage dendrogram. \textbf{(ii) Quasi-diagonalisation:} reorder assets
by the dendrogram's leaf order, concentrating large covariances near the
diagonal. \textbf{(iii) Recursive bisection:} split the ordered list in half,
allocate between halves in inverse proportion to their cluster variances, and
recurse. Because no matrix is inverted, HRP is robust to ill-conditioning and
beats mean--variance and inverse-variance portfolios out of sample
\cite{lopezdeprado2016hrp}.

Two structural observations drive this thesis. First, \emph{the dendrogram
exists only to produce an ordering}: allocation itself (step iii) consumes an
ordered list plus cluster variances, not the tree. Any procedure that yields
an economically meaningful asset ordering can replace the clustering step
wholesale --- the observation behind causal-ordered bisection
(Section~\ref{sec:d2}). Second, \emph{every input to steps (i)--(ii) is
symmetric by necessity}, not by accident: linkage clustering merges clusters
by comparing distances, and a distance is symmetric by definition, so the
pipeline is structurally blind to directionality. If the true data-generating
process contains one-way influence --- $A$ moves $B$ with a sign and
magnitude that $B$ does not reciprocate --- no correlation-fed hierarchical
allocator can represent that fact. This is the \emph{interface} deficit of
Section~\ref{sec:motivation}: not a flaw of correlation estimation but of
what the clustering step can accept. Making that interface accept directional
information, without breaking what makes HRP work, is challenge C2.

The limitation is family-wide, not HRP-specific. HRP seeded a family of
hierarchical allocators --- hierarchical clustering-based asset allocation
and its equal-risk-contribution variant \cite{raffinot2017hierarchical},
nested clustered optimisation \cite{lopezdeprado2020mlam} --- whose members
vary the linkage rule, the risk measure or the within-cluster optimiser,
but every one of which consumes a symmetric similarity structure at the
same first step. A decomposition of what directed structure is worth at
that interface therefore speaks to the family, not to one algorithm; HRP
is used here as its simplest and most-studied representative, so that any
effect found is attributable to the input rather than to allocator
elaborations.

The same author's later programme supplies the evaluation standard. The core
claim of the backtest-overfitting agenda \cite{bailey2014deflated,
lopezdeprado2023causal} is that the hard problem in quantitative finance is
not finding a configuration that looks good in a backtest, but establishing
that its performance is not an artefact of the search that found it. Two
instruments follow. The \emph{Probabilistic Sharpe Ratio} \cite{bailey2012psr}
replaces the Gaussian sampling theory of the Sharpe ratio --- untenable at the
excess kurtosis of ${\approx}\,\rsKurtosis$ measured here --- with a version
that corrects for the sample's skewness, kurtosis and length; the
\emph{Deflated Sharpe Ratio} \cite{bailey2014deflated} then raises the
benchmark from zero to the best Sharpe expected from $N$ independent trials
under the null of no skill, directly discounting selection bias. For
family-wise comparison against a benchmark, White's Reality Check
\cite{white2000reality} and Hansen's SPA test \cite{hansen2005spa} bootstrap
the null that \emph{no} candidate beats it, and the Model Confidence Set
\cite{hansen2011mcs} returns the subset of strategies statistically
indistinguishable from the best. This battery is applied to every claim in
Chapter~\ref{chap:evaluation}, over all \rsNtrials\ configurations this
project evaluated; several headline-looking effects do not survive it, and
reporting that is treated as part of the result rather than a failure.

\section{Theme II --- causal graphs on equities, without allocation:
Howard, Lohre and Mudde}
\label{sec:howard-theme}

Howard, Lohre and Mudde \cite{howard2025causal} is the closest prior work: a
large-scale application of DYNOTEARS to S\&P~500 constituents over
1989--2022 (${\approx}\,357$ names per estimated network), re-estimated on a
rolling basis --- structurally the same discovery machinery as
Chapter~\ref{chap:discovery}. Three of their empirical findings shape this
study's design. First, \emph{lagged edges are negligible} at daily-to-monthly
horizons: the contemporaneous block $W$ carries essentially all the
structure, which justifies this project's use of the contemporaneous
asset--asset block as ``the graph'' and its lag-1 truncation elsewhere.
Second, causal networks carry \emph{regime information that correlation
misses}: rising network density precedes lower returns and flags recessions,
and their centrality-based signal earns its alpha in stable rather than
stressed periods. This motivates the regime-conditional evaluation of
Section~\ref{sec:regimes} --- where, it turns out, the orientation effect
echoes their calm-regime pattern. Third, their overall verdict is measured
--- causal networks \emph{complement} rather than consistently beat
correlation-based methods --- a claim Q sharpens into a decomposition:
complement \emph{through what}, the skeleton or the orientations?

What they do \emph{not} do defines the gap this project fills. Their three
use-cases --- peer groups for factor neutralisation, a low-centrality factor,
and a density-based market-timing signal --- all consume the graph as a
similarity structure or a scalar diagnostic. The graph is never handed to an
allocator: no portfolio weights are ever computed \emph{from} the discovered
network, and the directionality of the edges is used only implicitly (via
centrality). Consequently their pipeline cannot ask, let alone answer, how
much of the graph's allocation value lies in its orientations. Tang
\cite{tang2024trading} builds trading signals from time-series causal
discovery (including VARLiNGAM) and finds computational cost the binding
constraint --- a warning this project met with a content-keyed discovery
cache (Section~\ref{sec:cache}) rather than by shrinking the experiment.

\section{Theme III --- the driver route and its instabilities:
Hierarchical Sensitivity Parity}
\label{sec:hsp-theme}

Hierarchical Sensitivity Parity \cite{rodriguez2023hsp} keeps HRP's
allocation but changes what assets are clustered on: instead of return
correlations, the similarity of assets' \emph{sensitivities to common
drivers}. For each asset a feed-forward neural network maps a set of
macro-financial drivers (rates, commodities, FX, volatility indices) to the
asset's return, and the sensitivity vector
$s_i = (\partial f_i/\partial d_1, \ldots, \partial f_i/\partial d_K)$,
extracted by automatic differentiation, embeds the asset in driver space; the
distance $\lVert s_i - s_j\rVert_2$ feeds the unchanged HRP machinery. The
theoretical appeal is the \emph{commonality principle}
\cite{rodriguez2023hsp,rodriguez2025causal}: if the drivers are genuinely the
common causes of the assets (in the sense of Reichenbach's common-cause
principle \cite{hitchcock2020reichenbach}), clustering in sensitivity space
diversifies systematic and idiosyncratic risk simultaneously.

Three weaknesses matter here, and the evaluation measures all three.
\emph{The selector}: the drivers are chosen by \emph{cumulative correlation}
with the assets --- the very quantity the causal motivation was meant to
transcend --- and the selection is demonstrably tautological in practice: on
this study's universe the correlation selector locks onto the same
international-equity ETFs and rates-curve points at essentially every
rebalance (Section~\ref{sec:r1}). \emph{The seeds}: the neural sensitivity
step makes the output seed-dependent --- re-running the identical pipeline
with \rsSeedN\ different network initialisations moves its net Sharpe across
a range ($\rsSeedMin$--$\rsSeedMax$) wider than the edge the method claims
over its own baseline, consistent with the concern that results reported for
such pipelines are sensitive to undisclosed seed choices
(Section~\ref{sec:seed-audit}). \emph{The loop}: the framework's natural
extension --- feeding realised performance back into driver selection ---
founders on a measurement problem: at a monthly rebalance the feedback reward
is a Sharpe ratio estimated from 21 daily observations, whose sampling error
exceeds the signal by more than an order of magnitude
(Section~\ref{sec:r1}). This project began as an attempt to repair the first
weakness, replacing correlational driver selection with causal discovery.
The evaluation shows that whole route to be dominated by allocating from the
asset--asset graph directly --- the result that reoriented the study to Q,
and the reason the driver route appears in this report as a documented dead
end rather than the destination.

\section{The causal-discovery toolkit}
\label{sec:cd-tools}

The pipeline needs three facts from the causal-discovery literature
\cite{spirtes2000causation,kaddour2022causalml}. \textbf{(i) Score-based
discovery scales and accepts priors.} Constraint-based methods (PC, FCI)
scale poorly with dimension and lean on faithfulness --- fragile in markets,
where offsetting mechanisms can leave causally-connected pairs with near-zero
marginal correlation \cite{howard2025causal}. NOTEARS \cite{zheng2018notears}
recasts acyclicity as the smooth constraint
$h(W)=\operatorname{tr}(e^{W\circ W})-d=0$, making DAG learning a continuous
optimisation; DYNOTEARS \cite{pamfil2020dynotears} extends it to the
structural VAR $X = XW + \sum_p Y_p A_p + Z$ with acyclicity on the
contemporaneous $W$ only, and natively supports forbidden-edge lists --- which
is how the directional prior of C1 is enforced. It is the primary method here.
\textbf{(ii) Non-Gaussianity identifies direction.} The LiNGAM family
\cite{shimizu2006lingam} exploits non-Gaussian residuals to identify a unique
causal ordering; VARLiNGAM \cite{hyvarinen2010varlingam} is its SVAR form and
serves as the robustness comparator --- its identification rests on an
assumption (heavy-tailed residuals) that financial returns satisfy at short
windows but progressively lose at longer ones as aggregation pulls residuals
toward Gaussianity, a mechanism the results repeatedly exhibit.
\textbf{(iii) A directed graph without trustworthy magnitudes exists in
closed form.} A ridge-regularised VAR(1) coefficient matrix is a
Granger-causal \cite{granger1969investigating} directed graph: it carries a
defensible directional \emph{ordering}, but its shrunken one-lag
coefficients are not structural effect sizes. That combination makes it a
natural control for dissociating what ``orientation'' contains --- the
arrows themselves versus magnitudes fit to propagate shocks. (A non-linear extension, NTS-NOTEARS \cite{sun2023ntsnotears},
was probed at reduced scale and left as future work: its prior-knowledge
mechanism transfers, but a full backtest costs ${\sim}10\times$ DYNOTEARS.)

\section{Research gap}
\label{sec:gap}

Table~\ref{tab:gap} positions the work; the \textsc{corr} control of
Section~\ref{sec:question} is exactly the HRP row. Reading across the
columns: HRP fixes the allocator but is blind to direction at the interface;
HSP re-routes information through driver sensitivities yet selects those
drivers by correlation and inherits neural instability; Howard et al.\
discover directed graphs at scale but never allocate from them. No prior work
sets hierarchical portfolio weights from a discovered asset--asset causal
graph; none allocates from an asymmetric matrix at all; and consequently none
can decompose the graph's allocation value into skeleton and orientation ---
the decomposition Q demands and
Chapters~\ref{chap:sensitivity}--\ref{chap:evaluation} deliver.

\begin{table}[t]
  \centering
  \small
  \caption{Positioning against the closest related approaches. The HRP row is
  the graph-blind \textsc{corr} control of this study.}
  \label{tab:gap}
  \begin{tabular}{l l c c c c}
    \toprule
    Approach & Clusters on & Causal & Weights & Asymmetric & Skeleton vs \\
             &             & discovery? & from graph? & allocation? & orientation? \\
    \midrule
    HRP \cite{lopezdeprado2016hrp} ($\equiv$ \textsc{corr}) & correlation & no & n/a & no & n/a \\
    HSP \cite{rodriguez2023hsp}            & driver sens.\    & no  & no  & no & no \\
    Howard et al.\ \cite{howard2025causal} & (no allocation)  & yes & no  & no & no \\
    \textbf{This work}                     & \textbf{causal graph} & \textbf{yes} & \textbf{yes} & \textbf{yes} & \textbf{yes} \\
    \bottomrule
  \end{tabular}
\end{table}
